\chapter{Двухшаговый моритов коннектор и граница смешанной кривизны}
\label{ch:v6-spectral-transition-morita-two-step-connector-gate}

\section{Последний внутренний кандидат}

Прямая стрелка от семейных полей $B,Q,T$ к слабому дублету отсутствует.
Тем не менее физический модуль одноформ
\begin{equation}
 \mathcal Y_\rho=\mathcal E_\rho\otimes\Lambda_{\rm ch}
\end{equation}
одновременно содержит семейный и хиггсовский множители. Поэтому остаётся
последний внутренний вопрос: может ли произведение двух уже существующих
одноформ породить во второй степени ненулевой портал
\begin{equation}
 \Tr(Q^2)H^\dagger H
 \label{eq:v6-two-step-target-portal}
\end{equation}
с фиксированными знаком и нормировкой?

Нулевая гипотеза утверждает, что смешанный член возникает из корректного
градуированного произведения или относительной моритовой кривизны.
Стоп-критерий --- одновременное выполнение трёх условий:
\begin{enumerate}
 \item смешанный оператор сокращается в градуированном квадрате;
 \item семейная бесследовая двухформа удаляется junk-идеалом;
 \item норма относительной кривизны распадается в сумму секторных норм.
\end{enumerate}

\section{Градуированный квадрат}

Пусть $D_F$ --- нечётный семейный оператор, $\gamma_F$ --- его
градуировка, а $D_H$ --- нечётный оператор хиггсовского ребра. Корректное
произведение имеет вид
\begin{equation}
 D_{\rm tot}=D_F\otimes1+\gamma_F\otimes D_H,
 \qquad
 \gamma_FD_F+D_F\gamma_F=0.
 \label{eq:v6-two-step-graded-product}
\end{equation}
Тогда
\begin{align}
 D_{\rm tot}^2
 &=D_F^2\otimes1+1\otimes D_H^2\nonumber\\
 &\quad+(D_F\gamma_F+\gamma_FD_F)\otimes D_H\nonumber\\
 &=D_F^2\otimes1+1\otimes D_H^2.
 \label{eq:v6-two-step-cross-cancellation}
\end{align}
Машинный остаток смешанного блока равен
$4.17\cdot10^{-16}$.

Если убрать $\gamma_F$, смешанный оператор действительно появляется;
контрольная норма равна $14.1797\ldots$. Но такая сумма не является
градуированным произведением нечётных циклов и не может использоваться
только ради получения желаемого портала. Формула
\eqref{eq:v6-two-step-cross-cancellation} является стандартной основой
произведения нечётных операторов в градуированной геометрии
\cite{two-step-kucerovsky1997,two-step-quillen1985}.

\begin{theorem}[Сокращение смешанного двухшагового оператора]
В корректном градуированном произведении семейной и наблюдаемой геометрий
смешанный член $D_F\otimes D_H$ тождественно сокращается. Поэтому сам
квадрат полного оператора не создаёт коннектор формы
\eqref{eq:v6-two-step-target-portal}.
\end{theorem}

\section{Обычные двухформы и junk-идеал}

Археология Тома IV уже вычислила полный представленный фактор
$\Omega_D^2$ семейной трёхузловой геометрии на общих фонах. На
физической половине ранг среднего блока после факторизации равен нулю:
\begin{equation}
 \rank\operatorname{pr}_{G,\rm particle}\Omega_D^2=0.
\end{equation}
На сопряжённой половине остаётся только центральная линия; ранг
бесследовой симметричной части также равен нулю.

Минимальная геометрия ранга один даёт ещё более прозрачную формулу:
\begin{equation}
 \pi(\Omega_u^2)=\mathbb C\rho\oplus QM_3Q,
 \qquad
 \pi(d\ker\pi_1)=QM_3Q,
\end{equation}
и потому
\begin{equation}
 \Omega_D^2\simeq\mathbb C\rho.
 \label{eq:v6-two-step-rank-one-quotient}
\end{equation}
Выживает радиальная опорная линия, но исчезает матричная форма
дополнительного угла. Следовательно, тензорное умножение выжившей линии
на $H^\dagger H$ способно повторить только радиальный канал, уже
опровергнутый предыдущим гейтом. Оно не восстанавливает $Q^2$.

Этот вывод использует стандартный представленный фактор форм Конна, а не
формальное умножение матриц до удаления вырожденных форм
\cite{two-step-schucker1996}.

\section{Относительная моритова кривизна}

Для переходного бимодуля $E=M_{20\times15}(\mathbb C)$ кривизна действует
как
\begin{equation}
 \mathcal R_E(\xi)=F_F\xi-\xi F_O.
\end{equation}
Её нормированный квадрат точно равен
\begin{equation}
 \tau_{300}(\mathcal R_E^2)
 =\tau_{20}(F_F^2)+\tau_{15}(F_O^2)
 -2\tau_{20}(F_F)\tau_{15}(F_O).
 \label{eq:v6-two-step-relative-curvature}
\end{equation}
Для бесследового параметра формы
\begin{equation}
 \tau_{20}(Q)=0,
\end{equation}
поэтому смешанный член исчезает. Остаётся
\begin{equation}
 \tau_{300}(\mathcal R_E^2)
 =\tau_{20}(F_F^2)+\tau_{15}(F_O^2).
 \label{eq:v6-two-step-additive-curvature}
\end{equation}
Численный остаток аддитивности равен $1.11\cdot10^{-16}$, а смешанная
производная по $\Tr Q^2$ и $|H|^2$ равна нулю.

Важно различать две формулы. Инвариант
\begin{equation}
 \tau_{20}(Q^2)\,H^\dagger H
\end{equation}
математически допустим как произведение двух секторных функционалов. Но
он не входит в квадрат относительной кривизны
\eqref{eq:v6-two-step-relative-curvature}. Его добавление потребовало бы
нового функционала высшей степени и нового правила фиксации коэффициента.

\begin{proposition}[Аддитивность вместо портала]
Центрированная относительная моритова кривизна фиксирует сумму семейной и
наблюдаемой норм, но не их произведение. Поэтому она не порождает
$\Tr(Q^2)H^\dagger H$.
\end{proposition}

\section{Проверка обычной суперсвязности}

Старый гейт градуированного соответствия уже разделил три кандидата:
\begin{align}
 Q_{\rm fam}^2\big|_{V_G}
 &=XX^\dagger+|\Phi|^2I_3,\nonumber\\
 d^2&=\Phi X,\nonumber\\
 [d,d^\dagger]-h
 &=XX^\dagger-|\Phi|^2I_3.
\end{align}
Первый даёт сумму матриц Грама, второй --- радиальную норму составного
пути $|\Phi|^2\Tr(X^\dagger X)$, а третий требует отдельно принятого
эрмитова примитива отображения момента. Ни один из первых двух не
порождает форму~\eqref{eq:v6-two-step-target-portal}; третий является
дополнительной структурой, а не следствием стандартной суперсвязности.

\begin{nogocriterion}[Запрет неградуированного или постфактум функционала]
Нельзя удалять $\gamma_F$ из произведения операторов, сохранять форму,
лежащую в junk-идеале, либо умножать две секторные нормы после вычисления
и объявлять полученный портал частью исходного родителя. Каждый такой шаг
меняет дифференциальное исчисление или функционал именно ради желаемого
смешанного члена.
\end{nogocriterion}

\section{Окончательная граница текущей архитектуры}

Последний внутренний кандидат не прошёл:
\begin{equation}
 \boxed{
 \text{градуированное сокращение}
 +\text{ junk формы}
 +\text{ аддитивность Мориты}
 \Longrightarrow \lambda_{QH}=0.
 }
\end{equation}
Это не запрещает портал во всех расширениях теории. Оно устанавливает,
что текущий набор полей, представленных форм, квадратичный функционал
кривизны и моритов контейнер его не выводят.

Следующий гейт
\begin{sloppypar}
\path{version6_spectral_transition_connector_architecture_branch_decision_gate}
\end{sloppypar}
должен зафиксировать развилку:
\begin{enumerate}
 \item сохранить спектральный переход и вихрь как два пока несвязанных
 эффективных сектора;
 \item либо объявить минимальное расширение новым бифундаментальным
 носителем или новым заранее заданным функционалом, после чего заново
 пройти представления, аномалии, нормировку и устойчивость.
\end{enumerate}

\begin{protocol}
\begin{enumerate}
 \item смешанный модуль одноформ: \textbf{сохранён};
 \item смешанный член градуированного квадрата:
 \textbf{тождественно сокращается};
 \item семейная форма в обычном $\Omega_D^2$:
 \textbf{удаляется junk-идеалом};
 \item фактор ранга один: \textbf{только $\mathbb C\rho$};
 \item центрированная относительная кривизна:
 \textbf{строго аддитивна};
 \item смешанная производная по $\Tr Q^2$ и $|H|^2$:
 \textbf{ноль};
 \item стандартная суперсвязность как источник портала:
 \textbf{не проходит};
 \item двухшаговый внутренний коннектор:
 \textbf{не получен};
 \item ненулевой портал текущим родителем:
 \textbf{не выведен};
 \item физическое замыкание: \textbf{не достигнуто}.
\end{enumerate}
\end{protocol}

Машинный сертификат сохранён в
\path{s2t_v6_spectral_transition_morita_two_step_connector_gate.py} и
\path{s2t_v6_spectral_transition_morita_two_step_connector_gate_results.json}.

\begin{thebibliography}{9}
\bibitem{two-step-kucerovsky1997}
D.~Kucerovsky,
\emph{The KK-Product of Unbounded Modules},
K-Theory 11 (1997), 17--34.

\bibitem{two-step-quillen1985}
D.~Quillen,
\emph{Superconnections and the Chern Character},
Topology 24 (1985), 89--95.

\bibitem{two-step-schucker1996}
T.~Sch\"ucker,
\emph{Geometries and Forces},
arXiv:hep-th/9605001.
\end{thebibliography}